\documentclass[11pt,a4paper]{article}
\usepackage[utf8]{inputenc}
\usepackage[T1]{fontenc}
\usepackage{lmodern}
\usepackage{microtype}
\usepackage[margin=2cm, top=2.5cm, bottom=2.5cm]{geometry}
\usepackage{fancyhdr}
\usepackage{titlesec}
\usepackage{amsmath, amssymb, amsthm}
\usepackage{mathtools}
\usepackage{bm}
\usepackage{booktabs}
\usepackage{array}
\usepackage{enumitem}
\usepackage[most]{tcolorbox}
\usepackage{xcolor}
\usepackage{hyperref}

% ------------------------------------------------------------------
% Palette
% ------------------------------------------------------------------
\definecolor{royalblue}{RGB}{28, 58, 185}
\definecolor{darkred}{RGB}{175, 22, 22}
\definecolor{burgundy}{RGB}{115, 8, 38}
\definecolor{bgblue}{RGB}{232, 238, 255}
\definecolor{bgred}{RGB}{255, 232, 232}
\definecolor{bgburgundy}{RGB}{255, 236, 242}
\definecolor{bggray}{RGB}{245, 245, 245}
\definecolor{framegray}{RGB}{190, 190, 190}

% Experiment colors (matching the dashboard CSS variables)
\definecolor{controlblue}{RGB}{100, 116, 139}
\definecolor{treatmentgreen}{RGB}{5, 150, 105}

\hypersetup{
  colorlinks=true,
  linkcolor=royalblue,
  urlcolor=royalblue,
  pdfauthor={Andr\'{e}s Ortega},
  pdftitle={A/B Testing: From Theory to Practice},
}

% ------------------------------------------------------------------
% Box environments
% ------------------------------------------------------------------
\newtcolorbox{formaldef}[1][Formal Definition]{
  enhanced, breakable,
  colback=bgblue,
  colframe=royalblue,
  fonttitle=\bfseries\color{royalblue},
  title={#1},
  titlerule=0.6pt,
  titlerule style={royalblue},
  sharp corners,
  boxrule=1.2pt,
  top=6pt, bottom=6pt, left=8pt, right=8pt,
}

\newtcolorbox{intuition}[1][Intuition]{
  enhanced, breakable,
  colback=bgred,
  colframe=darkred,
  fonttitle=\bfseries\color{darkred},
  title={#1},
  titlerule=0.6pt,
  titlerule style={darkred},
  sharp corners,
  boxrule=1.2pt,
  top=6pt, bottom=6pt, left=8pt, right=8pt,
}

\newtcolorbox{application}[1][Application]{
  enhanced, breakable,
  colback=bgburgundy,
  colframe=burgundy,
  fonttitle=\bfseries\color{burgundy},
  title={#1},
  titlerule=0.6pt,
  titlerule style={burgundy},
  sharp corners,
  boxrule=1.2pt,
  top=6pt, bottom=6pt, left=8pt, right=8pt,
}

\newtcolorbox{dimcheck}{
  enhanced, breakable,
  colback=bggray,
  colframe=framegray,
  fonttitle=\bfseries\color{black},
  title={Dimension Check},
  titlerule=0.6pt,
  sharp corners,
  boxrule=0.8pt,
  top=6pt, bottom=6pt, left=8pt, right=8pt,
}

% Dashboard link box
\newtcolorbox{dashlink}[1][Dashboard Connection]{
  enhanced, breakable,
  colback=bggray,
  colframe=framegray,
  fonttitle=\bfseries\color{black},
  title={#1},
  titlerule=0.6pt,
  sharp corners,
  boxrule=0.8pt,
  top=6pt, bottom=6pt, left=8pt, right=8pt,
}

% ------------------------------------------------------------------
% Section titles
% ------------------------------------------------------------------
\titleformat{\section}
  {\large\bfseries\color{darkred}}
  {\thesection.}{0.6em}{}
  [{\color{royalblue}\titlerule[0.8pt]}]

\titleformat{\subsection}
  {\normalsize\bfseries\color{royalblue}}
  {\thesubsection.}{0.5em}{}

% ------------------------------------------------------------------
% Header / Footer
% ------------------------------------------------------------------
\pagestyle{fancy}
\fancyhf{}
\fancyhead[L]{\small\color{darkred}\textbf{A/B Testing: From Theory to Practice}}
\fancyhead[R]{\small\color{royalblue}\thepage}
\renewcommand{\headrulewidth}{0.6pt}
\renewcommand{\headrule}{\color{royalblue}\hrule width\headwidth height\headrulewidth}

% ------------------------------------------------------------------
% Custom commands
% ------------------------------------------------------------------
\newcommand{\code}[1]{\texttt{#1}}
\newcommand{\scipy}[1]{\texttt{scipy.stats.#1}}
\newcommand{\numpy}[1]{\texttt{np.#1}}

\begin{document}

% ------------------------------------------------------------------
% Title block
% ------------------------------------------------------------------
\begin{center}
  {\LARGE\bfseries\color{darkred} A/B Testing: From Theory to Practice}\\[6pt]
  {\large\color{royalblue} A Bridge Document for the E-Commerce Landing Page Experiment}\\[8pt]
  {\small Andr\'{e}s Ortega \quad\textbullet\quad \today}
\end{center}
\vspace{4pt}
{\color{royalblue}\hrule height 1pt}
\vspace{10pt}

% ------------------------------------------------------------------
% Abstract
% ------------------------------------------------------------------
\begin{tcolorbox}[
  enhanced,
  colback=white,
  colframe=royalblue,
  boxrule=0.8pt,
  sharp corners,
  title={\bfseries\color{royalblue} Abstract},
  fonttitle=\bfseries,
  top=4pt, bottom=4pt
]
This document is a comprehensive guide to every statistical concept behind the A/B Test Lab dashboard.
It covers the full pipeline: from experimental design and randomization, through frequentist and Bayesian hypothesis testing,
to power analysis, sequential monitoring, and common pitfalls like Simpson's Paradox and multiple testing.
Each concept is presented in three layers: a formal mathematical definition, a geometric or analogy-based intuition,
and a concrete application grounded in the dashboard's actual backend code (\code{stats\_engine.py}).

\medskip
\textbf{Notation.}
$p_A, p_B$: true conversion rates for {\color{controlblue}control} and {\color{treatmentgreen}treatment}.
$\hat{p}_A, \hat{p}_B$: observed sample proportions.
$n_A, n_B$: sample sizes per group.
$\alpha$: significance level; $\beta$: Type~II error rate; $1-\beta$: power.
$\Phi$: standard normal CDF; $\Phi^{-1}$: its inverse.
$z_{\alpha/2} = \Phi^{-1}(1-\alpha/2)$, e.g.\ $z_{0.025} = 1.96$.
\end{tcolorbox}
\vspace{8pt}
\tableofcontents
\newpage

% ==================================================================
\section{Experimental Design Fundamentals}
% ==================================================================

\begin{formaldef}[Randomization and SRM Detection]
In a controlled experiment, each user $i$ is assigned to group $G_i \in \{A, B\}$ via:
\[
P(G_i = A) = \pi_A, \quad P(G_i = B) = 1 - \pi_A
\]
Typically $\pi_A = \pi_B = 0.5$. Randomization ensures that, in expectation, all covariates $X$ are balanced:
\[
E[X \mid G = A] = E[X \mid G = B]
\]
This is the foundation of causal inference: any observed difference in outcomes $Y$ can be attributed to the treatment, not to confounders.

\medskip
\textbf{Sample Ratio Mismatch (SRM).} Given expected counts $E_A = n \cdot \pi_A$ and $E_B = n \cdot \pi_B$, SRM is detected via chi-squared goodness-of-fit:
\[
\chi^2 = \frac{(n_A - E_A)^2}{E_A} + \frac{(n_B - E_B)^2}{E_B}, \quad \chi^2 \sim \chi^2_1 \text{ under } H_0
\]
Reject balance if $p < 0.01$. An SRM means the randomization mechanism is broken and all results become untrustworthy.
\end{formaldef}

\vspace{8pt}

\begin{intuition}[The Rigged Shuffle]
Think of randomization like shuffling a deck of cards and dealing into two piles. A fair shuffle balances \emph{everything}, even factors you have not thought of. That is its superpower.

SRM is the ``rigged shuffle'' detector. If you expected 50/50 and got 55/45 with 100K users, the chi-squared test will flag it: the shuffle was rigged. You should not trust any result from a rigged shuffle, even if the treatment looks amazing.

The \textbf{unit of randomization} matters: one user can generate many sessions. If you randomize by session, a power user who visits 50 times could see both variants, contaminating the experiment. Randomize by user: one user, one variant, always.
\end{intuition}

\vspace{8pt}

\begin{application}[Dashboard: Test Health Checks]
The \textbf{Executive Overview} tab shows the SRM test result under ``Test Health Checks.'' The backend function \code{sample\_ratio\_mismatch\_test(n\_control, n\_treatment)} computes the chi-squared statistic using \scipy{chi2.cdf(chi2, df=1)}, reports the p-value, and flags \code{is\_balanced = True} if $p > 0.01$.

\medskip
\textbf{Worked example from the dataset:} With $n_A = 145{,}310$ and $n_B = 145{,}274$ (total $290{,}584$):
\[
E_A = E_B = 145{,}292, \quad \chi^2 = \frac{(145310 - 145292)^2}{145292} + \frac{(145274 - 145292)^2}{145292} = \frac{324 + 324}{145292} \approx 0.0045
\]
$p \approx 0.95$. No SRM detected; the randomization is clean.
\end{application}

\vspace{8pt}

\begin{dashlink}
\textbf{Tab:} Executive Overview $\to$ Test Health Checks.\\
\textbf{Backend:} \code{routers/overview.py} calls \code{stats\_engine.sample\_ratio\_mismatch\_test()}.\\
\textbf{Scipy:} \code{scipy.stats.chi2.cdf(chi2, df=1)}.
\end{dashlink}


% ==================================================================
\section{Frequentist Inference for Proportions}
% ==================================================================

\subsection{The Two-Proportion Z-Test}

\begin{formaldef}[Z-Test for Two Proportions]
Let $X_A \sim \mathrm{Binomial}(n_A, p_A)$ and $X_B \sim \mathrm{Binomial}(n_B, p_B)$ be independent. Hypotheses:
\[
H_0: p_A = p_B \quad \text{vs.} \quad H_1: p_A \neq p_B
\]

\textbf{Pooled proportion} (under $H_0$):
\[
\hat{p} = \frac{X_A + X_B}{n_A + n_B}
\]

\textbf{Test statistic:}
\[
Z = \frac{\hat{p}_B - \hat{p}_A}{\underbrace{\sqrt{\hat{p}(1-\hat{p})\left(\dfrac{1}{n_A} + \dfrac{1}{n_B}\right)}}_{\text{pooled SE}}}
\]

Under $H_0$ and by the CLT, $Z \xrightarrow{d} N(0,1)$. Reject $H_0$ if $|Z| > z_{\alpha/2}$.

\medskip
\textbf{p-value:} $p = 2\,\Phi(-|Z|)$.

\medskip
\textbf{Confidence interval} (uses \emph{unpooled} SE, since we do not assume $p_A = p_B$ for the CI):
\[
\mathrm{SE}_{\mathrm{CI}} = \sqrt{\frac{\hat{p}_A(1-\hat{p}_A)}{n_A} + \frac{\hat{p}_B(1-\hat{p}_B)}{n_B}}, \qquad \mathrm{CI}_{1-\alpha}: (\hat{p}_B - \hat{p}_A) \pm z_{\alpha/2} \cdot \mathrm{SE}_{\mathrm{CI}}
\]

\textbf{Why pooled for the test but unpooled for the CI?} Under $H_0$, the best estimate of the common $p$ is the pooled proportion. But the CI estimates the \emph{difference} $p_B - p_A$ without assuming equality, so each variance is estimated separately.
\end{formaldef}

\vspace{8pt}

\begin{intuition}[Signal-to-Noise Ratio]
The numerator $\hat{p}_B - \hat{p}_A$ is the \textbf{signal}: the observed difference. The denominator (SE) is the \textbf{noise}: how much random variation we would expect from chance alone. $Z$ measures how many ``noise units'' the signal spans. If $|Z| > 2$, the signal is louder than what chance typically produces.

\medskip
\textbf{The CI tells you the range of plausible truths.} If the 95\% CI for the difference is $(0.001, 0.008)$, the true lift is somewhere in there with 95\% coverage. If the CI includes zero, the data is consistent with no effect.
\end{intuition}

\vspace{8pt}

\begin{application}[Dashboard: Frequentist Tab]
The backend function \code{z\_test\_proportions(n\_a, conv\_a, n\_b, conv\_b)} returns \code{\{z\_stat, p\_value, ci\_lower, ci\_upper, diff\}}. It uses \scipy{norm.cdf(abs(z\_stat))} for the p-value and \scipy{norm.ppf(0.975)} for the critical $z$ value.

The \textbf{Frequentist Analysis} tab displays:
\begin{itemize}[nosep]
  \item Confidence interval visualization (horizontal bar showing control CI, treatment CI, and overlap)
  \item P-value context: where it falls relative to $\alpha = 0.05$
  \item The 2$\times$2 contingency table with observed vs.\ expected counts
\end{itemize}
\end{application}

\subsection{Chi-Squared Test of Independence}

\begin{formaldef}[Chi-Squared for a 2$\times$2 Table]
For the contingency table with observed counts $O_{ij}$ and expected counts $E_{ij} = R_i \cdot C_j / N$:
\[
\chi^2 = \sum_{i,j} \frac{(O_{ij} - E_{ij})^2}{E_{ij}}, \qquad \chi^2 \sim \chi^2_1
\]

\textbf{Equivalence:} For a 2$\times$2 table, $\chi^2 = Z^2$ exactly. They always give the same p-value for a two-sided test. The chi-squared test is the z-test in disguise; squaring removes the direction.

\medskip
\textbf{Cram\'{e}r's V} (effect size for contingency tables):
\[
V = \sqrt{\frac{\chi^2}{n}}
\]
\end{formaldef}

\vspace{8pt}

\begin{intuition}[Z-Test in Disguise]
Squaring the z-statistic removes the sign (direction), which is why chi-squared is inherently two-sided. If you need a one-sided test, use the z-test directly. For A/B testing, we almost always want to know the direction (``is B \emph{better}?''), so the z-test is more informative. The chi-squared is a redundant check that should give the same answer.
\end{intuition}

\vspace{8pt}

\begin{application}[Backend Implementation]
\code{chi\_squared\_test(n\_a, conv\_a, n\_b, conv\_b)} builds the 2$\times$2 contingency table and calls \scipy{chi2\_contingency(table)}, returning $\chi^2$, p-value, and Cram\'{e}r's V. The Frequentist tab displays all three alongside the z-test results.
\end{application}

\subsection{Wilson Score Confidence Interval}

\begin{formaldef}[Wilson Score Interval]
The standard Wald interval $\hat{p} \pm z\sqrt{\hat{p}(1-\hat{p})/n}$ can produce intervals outside $[0,1]$ and has zero width when $\hat{p} \in \{0, 1\}$.

The Wilson interval is derived by inverting the z-test. Find all $p_0$ such that:
\[
\left|\frac{\hat{p} - p_0}{\sqrt{p_0(1-p_0)/n}}\right| \leq z_{\alpha/2}
\]

Squaring and solving the resulting quadratic:
\[
\tilde{p} = \frac{\hat{p} + z^2/(2n)}{1 + z^2/n}, \qquad \text{Wilson CI} = \tilde{p} \pm \frac{z}{1 + z^2/n}\sqrt{\frac{\hat{p}(1-\hat{p})}{n} + \frac{z^2}{4n^2}}
\]

The center $\tilde{p}$ is a \emph{shrunken} estimate that pulls $\hat{p}$ toward $0.5$, more so for small $n$.
\end{formaldef}

\vspace{8pt}

\begin{intuition}[Pulling Toward Ignorance]
The Wald interval says ``I trust my data completely.'' The Wilson interval says ``I trust my data, but I hedge toward 50\% a little, especially when I have few observations.'' This prevents the embarrassing case of observing 0 failures in 10 trials and concluding the failure rate is exactly zero.
\end{intuition}

\vspace{8pt}

\begin{application}[Dashboard: Wilson CIs per Group]
\code{wilson\_confidence\_interval(successes, trials, alpha)} uses \scipy{norm.ppf(1 - alpha/2)} and returns the (lower, upper) bounds. The Frequentist tab calls this separately for control and treatment to display per-group CIs, which is more informative than just the CI on the difference.
\end{application}

\subsection{Cohen's h Effect Size}

\begin{formaldef}[Cohen's h for Proportions]
The arcsine transformation stabilizes variance:
\[
h = 2\arcsin\!\left(\sqrt{p_B}\right) - 2\arcsin\!\left(\sqrt{p_A}\right)
\]

\textbf{Why arcsine?} The transformation $\phi = 2\arcsin(\sqrt{p})$ is variance-stabilizing: $\mathrm{Var}(\hat{p}) = p(1-p)/n$ depends on $p$, but $\mathrm{Var}(\phi) \approx 1/n$, independent of $p$.

\medskip
\begin{tabular}{@{}ll@{}}
\toprule
$|h|$ & Interpretation \\
\midrule
$< 0.20$ & Small effect \\
$0.20$--$0.50$ & Medium effect \\
$\geq 0.50$ & Large effect \\
\bottomrule
\end{tabular}

\medskip
\textbf{Example:} $p_A = 0.10$, $p_B = 0.12$ (20\% relative lift). $h = 2\arcsin(\sqrt{0.12}) - 2\arcsin(\sqrt{0.10}) = 0.064$, a very small effect. Large relative lifts on small baselines are hard to detect because the absolute effect in standardized units is tiny.
\end{formaldef}

\vspace{8pt}

\begin{intuition}[Proportions Near 0 or 1 Are Harder]
Imagine stretching a rubber band between 0 and 1. The arcsine transformation stretches the ends (near 0 and 1) and compresses the middle (near 0.5). A change from 0.01 to 0.02 maps to a larger $h$ than you might expect, while 0.49 to 0.50 maps to a smaller $h$. This reflects reality: detecting changes near the extremes is genuinely harder because $p(1-p)$ is smallest there.
\end{intuition}

\vspace{8pt}

\begin{application}[Effect Size Gauge on the Dashboard]
\code{cohens\_h(p1, p2)} computes $h$ and returns the interpretation string. The Frequentist tab displays this as an ``Effect Size Gauge'' with small/medium/large labels, giving stakeholders a quick read on whether the detected difference is practically meaningful, not just statistically significant.
\end{application}

\subsection{Proper p-Value Interpretation}

\begin{formaldef}[What a p-Value Is and Is Not]
Formally:
\[
p\text{-value} = P(\text{data as extreme or more extreme} \mid H_0 \text{ is true})
\]

\textbf{Common misinterpretations:}
\begin{enumerate}[nosep]
  \item ``There is a 5\% probability the null is true.'' \textbf{Wrong.} The p-value conditions on $H_0$; computing $P(H_0 \mid \text{data})$ requires Bayes' theorem and a prior.
  \item ``The effect size is $1 - p$.'' \textbf{Wrong.} p-values say nothing about effect magnitude.
  \item ``$p = 0.049$ is meaningfully different from $p = 0.051$.'' \textbf{Wrong.} The $\alpha = 0.05$ threshold is a convention, not a physical constant.
  \item ``Failing to reject $H_0$ means $H_0$ is true.'' \textbf{Wrong.} The test may be underpowered (see Section~\ref{sec:power}).
\end{enumerate}
\end{formaldef}

\vspace{8pt}

\begin{intuition}[The Courtroom Analogy]
A p-value is like the prosecution's evidence in a trial where the defendant ($H_0$) is presumed innocent. A small p-value means the evidence is strong, but even overwhelming evidence does not tell you the \emph{probability} of innocence. The jury can say ``guilty beyond reasonable doubt'' (reject $H_0$) or ``not proven'' (fail to reject), but never ``there is a 3\% chance the defendant is innocent.''
\end{intuition}

\vspace{8pt}

\begin{application}[Dashboard: P-Value Context]
The Frequentist tab shows the p-value not just as a number but as a visual showing where it falls relative to $\alpha$, preventing the common mistake of treating $p = 0.049$ and $p = 0.051$ as fundamentally different outcomes.
\end{application}

\vspace{8pt}

\begin{dashlink}
\textbf{Tab:} Frequentist Analysis.\\
\textbf{Backend:} \code{routers/frequentist.py} calls \code{z\_test\_proportions()}, \code{chi\_squared\_test()}, \code{wilson\_confidence\_interval()} (for both groups), and \code{cohens\_h()}.\\
\textbf{Additional:} Welch's t-test via \scipy{ttest\_ind(equal\_var=False)} for continuous metrics (revenue, session duration).
\end{dashlink}


% ==================================================================
\section{Bayesian Inference for A/B Testing}
% ==================================================================

\subsection{The Beta-Binomial Conjugate Model}

\begin{formaldef}[Beta Prior and Posterior Update]
\textbf{Prior:}
\[
p_A \sim \mathrm{Beta}(\alpha_A, \beta_A), \quad f(p; \alpha, \beta) = \frac{p^{\alpha-1}(1-p)^{\beta-1}}{B(\alpha, \beta)}
\]

\textbf{Common prior choices:}
\begin{itemize}[nosep]
  \item $\mathrm{Beta}(1, 1)$: uniform (non-informative). The dashboard uses this.
  \item $\mathrm{Beta}(0.5, 0.5)$: Jeffreys prior (invariant under reparameterization).
  \item $\mathrm{Beta}(100, 9900)$: informative, encoding belief that $p \approx 0.01$.
\end{itemize}

\textbf{Likelihood:} $X_A \mid p_A \sim \mathrm{Binomial}(n_A, p_A)$.

\textbf{Posterior} (conjugacy):
\[
p_A \mid X_A \sim \mathrm{Beta}(\alpha_A + X_A, \;\beta_A + n_A - X_A)
\]

The update is simply \emph{addition}: add successes to $\alpha$, add failures to $\beta$. The posterior mean is a weighted average between the prior mean and the MLE:
\[
E[p_A \mid X_A] = \frac{\alpha_A + X_A}{\alpha_A + \beta_A + n_A}
\]
\end{formaldef}

\vspace{8pt}

\begin{intuition}[The Belief Histogram and Imaginary Data]
The Beta distribution is a ``belief histogram'' over where the true conversion rate lies. $\mathrm{Beta}(1,1)$ is flat: equal belief everywhere. After observing 100 conversions and 9,900 non-conversions, your belief becomes $\mathrm{Beta}(101, 9901)$: a sharp bump centered at ${\sim}0.01$.

\medskip
\textbf{The prior as imaginary data.} $\mathrm{Beta}(\alpha, \beta)$ is equivalent to having already seen $\alpha - 1$ successes and $\beta - 1$ failures. $\mathrm{Beta}(1,1)$ means zero imaginary observations. The more imaginary data you inject, the harder it is for real data to move your beliefs.

\medskip
\textbf{The see-saw.} The prior is one end, the data is the other. With little data, the see-saw tips toward the prior. With lots of data, it tips toward the observed rate. The fulcrum position is $n_{\mathrm{prior}} / (n_{\mathrm{prior}} + n_{\mathrm{data}})$ where $n_{\mathrm{prior}} = \alpha + \beta$.
\end{intuition}

\vspace{8pt}

\begin{application}[Dashboard: Posterior Distribution Chart]
\code{beta\_posterior(prior\_alpha, prior\_beta, successes, failures)} returns the updated $(\alpha, \beta)$. The backend uses $\mathrm{Beta}(1, 1)$ as the prior (constants \code{PRIOR\_ALPHA = 1.0}, \code{PRIOR\_BETA = 1.0}).

\code{posterior\_pdf(alpha, beta, x\_points=200)} evaluates the density via \scipy{beta.pdf(x, alpha, beta)} at 200 points, sent to the frontend as \code{\{x: [...], pdf: [...]\}} for the overlapping area chart on the Bayesian tab.
\end{application}

\subsection{P(B $>$ A) via Monte Carlo}

\begin{formaldef}[Probability That Treatment Beats Control]
Draw $S$ samples from each posterior:
\[
p_A^{(s)} \sim \mathrm{Beta}(\alpha_A', \beta_A'), \quad p_B^{(s)} \sim \mathrm{Beta}(\alpha_B', \beta_B'), \quad s = 1, \ldots, S
\]
Then:
\[
\widehat{P}(p_B > p_A \mid \text{data}) = \frac{1}{S}\sum_{s=1}^{S} \mathbf{1}\!\left[p_B^{(s)} > p_A^{(s)}\right]
\]

With $S = 100{,}000$, the Monte Carlo standard error is $\leq \sqrt{0.25/100000} \approx 0.0016$.

\medskip
\textbf{Closed-form} (Evan Miller): For equal priors, there is an exact formula using the regularized incomplete beta function, but it is computationally expensive for large parameters. Monte Carlo is more practical.
\end{formaldef}

\vspace{8pt}

\begin{intuition}[The Question Everyone Actually Wants Answered]
When a product manager asks ``is B better?'', they want a probability. The frequentist answers: ``if there were truly no difference, we would see data this extreme only 3\% of the time.'' That answers a question nobody asked.

The Bayesian $P(B > A) = 0.97$ is direct: ``there is a 97\% probability that B has a higher conversion rate.'' This is the natural language of decision-making.
\end{intuition}

\vspace{8pt}

\begin{application}[Dashboard: Probability Donut Chart]
\code{probability\_b\_beats\_a(post\_a, post\_b, n\_simulations=100000)} uses \numpy{random.default\_rng(42).beta()} for reproducibility. The Bayesian tab displays this as a donut chart labeled ``Probability of Being Best.''
\end{application}

\subsection{Expected Loss}

\begin{formaldef}[Expected Loss (Regret)]
The loss of choosing B when the true parameters are $(p_A, p_B)$:
\[
L(\text{choose } B; p_A, p_B) = \max(p_A - p_B, \; 0)
\]
If B is worse, we lose the difference; if B is better, we lose nothing.

\textbf{Expected loss} via Monte Carlo:
\[
\widehat{E}\bigl[\text{Loss}(\text{choose } B)\bigr] = \frac{1}{S}\sum_{s=1}^{S} \max\!\left(p_A^{(s)} - p_B^{(s)}, \; 0\right)
\]

\textbf{Decision rule:} Ship B when $E[\text{Loss}(\text{choose } B)] < \varepsilon$, where $\varepsilon$ is a business-defined threshold. Unlike the frequentist $\alpha$, $\varepsilon$ has a direct business interpretation: ``we accept a maximum average loss of $\varepsilon$ percentage points of conversion rate.''
\end{formaldef}

\vspace{8pt}

\begin{intuition}[Insurance Pricing]
Think of choosing variant B as buying insurance. The expected loss is the premium: the average amount you would lose if B turns out worse, weighted by how likely that is. If the premium is negligibly small (below $\varepsilon$), the ``insurance'' is practically free: go with B.
\end{intuition}

\vspace{8pt}

\begin{application}[Dashboard: Expected Loss Display]
\code{expected\_loss(post\_a, post\_b, n\_simulations=100000)} uses \numpy{maximum(diff, 0)} for the one-sided max and returns \code{\{loss\_a, loss\_b\}}. The Bayesian tab shows both losses, allowing the viewer to see not just which variant is better, but how much they would lose by picking the wrong one.
\end{application}

\subsection{Credible Intervals}

\begin{formaldef}[Equal-Tailed CI vs.\ Highest Density Interval]
\textbf{Equal-tailed credible interval (ETI):} $[q_{\alpha/2}, q_{1-\alpha/2}]$ where $q_p$ is the $p$-th quantile of the posterior. Excludes $\alpha/2$ probability from each tail.

\textbf{Highest Density Interval (HDI):} The narrowest interval containing $1-\alpha$ probability mass. For any point inside the HDI, the posterior density is higher than for any point outside.

\medskip
For symmetric posteriors (large $n$, moderate $p$): ETI $\approx$ HDI. For skewed posteriors (small $n$, extreme $p$): HDI is narrower and more informative.
\end{formaldef}

\vspace{8pt}

\begin{application}[Backend: Beta Quantiles]
\code{credible\_interval(alpha, beta, level=0.95)} uses \scipy{beta.ppf()} to compute the equal-tailed interval. The dashboard shows side-by-side 95\% credible interval bars for control and treatment on the Bayesian tab.
\end{application}

\vspace{8pt}

\begin{dashlink}
\textbf{Tab:} Bayesian Analysis.\\
\textbf{Backend:} \code{routers/bayesian.py} calls all four functions above, plus builds a comparison table against the frequentist results. Agreement check: frequentist significant $\approx$ Bayesian decisive ($P(B>A) > 0.95$ or $< 0.05$).
\end{dashlink}


% ==================================================================
\section{Frequentist vs.\ Bayesian: A Comparison}
% ==================================================================

\begin{formaldef}[Side-by-Side Comparison]
\begin{tabular}{@{}p{3cm}p{5.5cm}p{5.5cm}@{}}
\toprule
\textbf{Dimension} & \textbf{Frequentist} & \textbf{Bayesian} \\
\midrule
Probability & Long-run frequency & Degree of belief \\
Parameters & Fixed, unknown constants & Random variables with distributions \\
Output & p-values, confidence intervals & Posteriors, credible intervals \\
Error control & Type I ($\alpha$), Type II ($\beta$) rates & Expected loss, posterior risk \\
Sample size & Must be fixed before experiment & Can be flexible (with caveats) \\
Prior info & Not formally incorporated & Encoded via prior distributions \\
Interval meaning & ``95\% of CIs from repeated experiments contain the truth'' & ``95\% probability the truth is in this interval'' \\
\bottomrule
\end{tabular}

\medskip
\textbf{Bernstein--von Mises theorem:} With a flat prior and large samples, Bayesian credible intervals converge to frequentist confidence intervals:
\[
p \mid \text{data} \xrightarrow{d} N\!\left(\hat{p}_{\mathrm{MLE}}, \; \frac{1}{I(\hat{p}_{\mathrm{MLE}})}\right)
\]
In finite samples with informative priors, the approaches can diverge meaningfully.
\end{formaldef}

\vspace{8pt}

\begin{intuition}[The Weather Forecast]
\textbf{Frequentist:} ``Under long-run repetitions of similar atmospheric conditions, it rains 30\% of the time.'' Technically correct, but awkward.

\textbf{Bayesian:} ``There is a 30\% chance of rain tomorrow.'' Natural, actionable, what people actually want.

\medskip
\textbf{When frequentist is better:} Regulatory contexts (FDA, financial audits) where pre-specified error rates matter; when there is no reasonable prior; when transparency is paramount.

\textbf{When Bayesian is better:} Decision-making under uncertainty; small samples where the prior regularizes; when stakeholders want probabilities, not p-values; sequential decisions.
\end{intuition}

\vspace{8pt}

\begin{application}[Dashboard: Comparison Table]
The Bayesian tab's comparison table (built in \code{routers/bayesian.py}) shows both frameworks' conclusions side by side: point estimate, interval, and conclusion. When they agree, the decision is easy. When they disagree (e.g., $p = 0.049$ but $P(B>A) = 0.83$), it signals that evidence is weak and more data is needed.
\end{application}


% ==================================================================
\section{Power Analysis}
\label{sec:power}
% ==================================================================

\subsection{Sample Size for Two Proportions}

\begin{formaldef}[Sample Size Formula]
For a two-sided z-test comparing $p_A$ vs.\ $p_B = p_A + \delta$ with significance $\alpha$ and power $1-\beta$:
\[
n = \frac{\left(z_{\alpha/2}\sqrt{2\bar{p}(1-\bar{p})} + z_{\beta}\sqrt{p_A(1-p_A) + p_B(1-p_B)}\right)^2}{(p_B - p_A)^2}
\]
where $\bar{p} = (p_A + p_B)/2$.

\medskip
\textbf{Simplified approximation} (when $p_B \approx p_A$):
\[
n \approx \frac{(z_{\alpha/2} + z_\beta)^2 \cdot 2p(1-p)}{(p_B - p_A)^2}
\]

This gives $n$ per group. Total users $= 2n$.
\end{formaldef}

\vspace{8pt}

\begin{intuition}[Trip Planning]
Power analysis is trip planning. ``Do I have enough gas (users) to reach my destination (detect the effect)?''

\begin{itemize}[nosep]
  \item $n$ = tank size
  \item MDE = distance to destination
  \item $\alpha$ = speed limit (strictness of the rules)
  \item $1-\beta$ = probability of arriving (80\% power = 80\% chance of ``getting there'')
\end{itemize}

\textbf{The $1/\sqrt{n}$ law.} Halving the MDE requires \emph{quadrupling} the sample size. Standard errors shrink as $1/\sqrt{n}$, not $1/n$. Going from 1\% MDE to 0.5\% MDE means going from 10K to 40K users per group. This ``law of diminishing returns'' is the fundamental constraint of experimentation.
\end{intuition}

\vspace{8pt}

\begin{application}[Dashboard: Power Calculator]
The Power \& Design tab has interactive sliders for baseline rate, MDE, $\alpha$, and power. The backend \code{required\_sample\_size(baseline\_rate, mde, alpha, power)} uses \scipy{norm.ppf(1 - alpha/2)} and \scipy{norm.ppf(power)} to compute the formula above.

\medskip
\textbf{Worked example:} $p_A = 0.12$, MDE $= 0.01$, $\alpha = 0.05$, power $= 0.80$:
\[
z_\alpha = 1.96, \quad z_\beta = 0.84, \quad p_B = 0.13
\]
\[
n = \frac{(1.96\sqrt{2 \cdot 0.125 \cdot 0.875} + 0.84\sqrt{0.12 \cdot 0.88 + 0.13 \cdot 0.87})^2}{(0.01)^2} \approx 14{,}752 \text{ per group}
\]
\end{application}

\subsection{Minimum Detectable Effect and Power Curves}

\begin{formaldef}[MDE and Power Function]
\textbf{MDE} at a given $n$: the smallest $\delta$ such that $n \geq n_{\mathrm{required}}(\delta)$. Found by binary search inverting the sample size formula.

\medskip
\textbf{Power as a function of effect size $\delta$:}
\[
\mathrm{Power}(\delta) = \Phi\!\left(\frac{|\delta|\sqrt{n}}{\sqrt{p_A(1-p_A) + p_B(1-p_B)}} - z_{\alpha/2}\right)
\]

Key properties: Power$(0) = \alpha$; power increases with $|\delta|$ and $n$; decreases with $\alpha$.
\end{formaldef}

\vspace{8pt}

\begin{intuition}[The Sensitivity Curve]
A power curve is a ``sensitivity diagram'': it shows how likely your experiment is to detect effects of various sizes. The curve always starts at $\alpha$ (at zero effect, ``detection'' is just a false positive) and rises toward 1. The MDE is the effect size where the curve crosses 0.80.

\textbf{Why 80\% and not 95\%?} Going to 95\% power roughly doubles sample size. Most product teams prefer speed, accepting that 1 in 5 real effects will be missed.
\end{intuition}

\vspace{8pt}

\begin{application}[Dashboard: MDE Curve and Runtime]
\code{minimum\_detectable\_effect(n, baseline\_rate)} uses binary search. \code{power\_curve(baseline\_rate, n)} evaluates power at 25 predefined effect sizes using \scipy{norm.cdf()}. The tab also estimates runtime:
\[
\text{Days} = \frac{n_{\mathrm{required}} - n_{\mathrm{current}}}{\text{daily\_traffic} / 2}
\]
Daily traffic is estimated as the mean daily observation count from the data.
\end{application}

\vspace{8pt}

\begin{dashlink}
\textbf{Tab:} Power \& Design.\\
\textbf{Backend:} \code{routers/power.py} calls \code{required\_sample\_size()}, \code{minimum\_detectable\_effect()}, and \code{power\_curve()}.\\
\textbf{Scipy:} \code{scipy.stats.norm.ppf()}, \code{scipy.stats.norm.cdf()}.
\end{dashlink}


% ==================================================================
\section{Sequential Testing}
\label{sec:sequential}
% ==================================================================

\subsection{Why Peeking Inflates Type I Error}

\begin{formaldef}[The Peeking Problem]
If you peek at data $k$ times during collection and stop the first time $|Z| > z_{\alpha/2}$, the actual Type~I error rate under $H_0$ is:
\[
\alpha_{\mathrm{inflated}} \approx 1 - (1 - \alpha)^k
\]
For $k = 5$: $\alpha_{\mathrm{inflated}} \approx 1 - 0.95^5 = 0.226$ (22.6\%, not 5\%).

\medskip
\textbf{The mathematical core.} Under $H_0$, the cumulative z-statistic at information fraction $t$ follows approximately $Z(t) \approx W(t)/\sqrt{t}$ where $W(t)$ is a standard Brownian motion. Since Brownian motion is recurrent, it will \emph{eventually} cross any finite boundary. Continuous peeking with fixed boundaries guarantees a false positive with probability~1.
\end{formaldef}

\vspace{8pt}

\begin{intuition}[Lottery Tickets]
Each time you check a running experiment, you buy a lottery ticket. Even though each ticket has only a 5\% chance of ``winning'' (false positive), buying 10 tickets gives you a 40\% chance. Sequential testing does not stop you from buying tickets; it makes each ticket progressively harder to win, so the total chance stays at 5\%.

\medskip
\textbf{The drunk on a number line.} Under the null, your cumulative test statistic wanders like a drunk person. Given enough time, the drunk will wander arbitrarily far from the origin. Fixed boundaries ($\pm 1.96$) will \emph{always} be crossed eventually. Sequential testing uses expanding boundaries to account for this inevitable wandering.
\end{intuition}

\vspace{8pt}

\begin{application}[Dashboard: Daily P-Value Evolution]
The Sequential tab shows how the p-value fluctuates over the 22-day experiment period. This visualization directly demonstrates the peeking problem: some days the p-value dips below 0.05, then bounces back. A naive analyst who stopped on those days would have declared a false winner.
\end{application}

\subsection{O'Brien--Fleming Boundaries}

\begin{formaldef}[Group Sequential Design]
Divide the experiment into $K$ stages at information fractions $t_1 < \cdots < t_K = 1$. At each stage $k$, compare $|Z_k|$ to boundary $c_k$:
\begin{itemize}[nosep]
  \item Stop and reject $H_0$ if $|Z_k| \geq c_k$
  \item Continue if $|Z_k| < c_k$ for $k < K$
\end{itemize}

Boundaries $\{c_k\}$ are chosen so the overall Type~I error is exactly $\alpha$:
\[
P\!\left(\bigcup_{k=1}^{K}\{|Z_k| \geq c_k\} \;\middle|\; H_0\right) = \alpha
\]

\textbf{O'Brien--Fleming boundary:}
\[
c_k^{\mathrm{OBF}} = c \cdot \sqrt{K/k}
\]
where $c$ is calibrated for overall $\alpha$.

\medskip
\begin{tabular}{@{}cccc@{}}
\toprule
Stage $k$ & Info fraction & OBF boundary & Approx.\ $\alpha$ spent \\
\midrule
1 & 0.20 & 4.56 & 0.000005 \\
2 & 0.40 & 3.23 & 0.0013 \\
3 & 0.60 & 2.63 & 0.0085 \\
4 & 0.80 & 2.28 & 0.0228 \\
5 & 1.00 & 2.04 & 0.0417 \\
\bottomrule
\end{tabular}

\medskip
Key: very conservative early (overwhelming evidence required), final boundary (2.04) close to fixed-sample (1.96). Total sample size increase: only ${\sim}3\%$ more.

\medskip
\textbf{Alpha spending function} (Lan--DeMets OBF-type):
\[
\alpha^*_{\mathrm{OBF}}(t) = 2 - 2\Phi\!\left(\frac{z_{\alpha/2}}{\sqrt{t}}\right)
\]
This does not require pre-specifying the number or timing of interim analyses.
\end{formaldef}

\vspace{8pt}

\begin{intuition}[The Skeptical Judge]
Early in a trial, the OBF boundary is extremely high, like a judge who says ``extraordinary claims require extraordinary evidence.'' As more data accumulates, the judge relaxes. This is sensible because early results are noisy (small $n$), so we should demand stronger evidence.

\medskip
\textbf{Alpha spending as a budget.} You have \$0.05 of ``significance dollars.'' The spending function decides how to spread the budget over time. OBF saves most for the end. You can never overspend; once the budget is exhausted, you stop.
\end{intuition}

\vspace{8pt}

\begin{application}[Dashboard: Sequential Chart]
The backend \code{obrien\_fleming\_bounds(n\_looks, alpha)} uses \scipy{norm.ppf(1 - alpha/2)} and computes $c_k = z_\alpha / \sqrt{k/K}$ for each day of the experiment. \code{cumulative\_stats\_by\_day(df)} computes running z-statistics by calling \code{z\_test\_proportions()} on cumulative data each day.

The Sequential tab overlays the z-statistic trajectory on the OBF boundaries, and highlights the \textbf{optimal stopping point}: the first day where $|z| \geq c_k$, if any.
\end{application}

\vspace{8pt}

\begin{dashlink}
\textbf{Tab:} Sequential Monitoring.\\
\textbf{Backend:} \code{routers/sequential.py} calls \code{cumulative\_stats\_by\_day()} and \code{obrien\_fleming\_bounds()}.\\
\textbf{Output:} cumulative stats, boundaries, daily p-values, and optimal stopping point.
\end{dashlink}


% ==================================================================
\section{Simpson's Paradox}
\label{sec:simpson}
% ==================================================================

\begin{formaldef}[When Aggregation Reverses the Truth]
Let $X$ be treatment, $Y$ outcome, $Z$ a confounding variable. Simpson's Paradox occurs when:
\[
P(Y{=}1 \mid X{=}B, Z{=}z) > P(Y{=}1 \mid X{=}A, Z{=}z) \quad \forall\, z
\]
but
\[
P(Y{=}1 \mid X{=}B) < P(Y{=}1 \mid X{=}A)
\]
B beats A in every subgroup, but A beats B in aggregate.

\medskip
\textbf{How it happens.} The marginal rate is a weighted average:
\[
p_X = \sum_z w_{X,z} \cdot p_{X,z}, \quad w_{X,z} = P(Z{=}z \mid X)
\]
When the \textbf{weights} $w_{X,z}$ differ between treatments and correlate with the subgroup rates, aggregation distorts the comparison.

\medskip
\textbf{Numerical example:}

\medskip
\begin{tabular}{@{}lccc@{}}
\toprule
 & Subgroup 1 (easy) & Subgroup 2 (hard) & Aggregate \\
\midrule
Treatment A & 81/87 = 93.1\% & 192/263 = 73.0\% & 273/350 = 78.0\% \\
Treatment B & 234/270 = 86.7\% & 55/80 = 68.8\% & 289/350 = 82.6\% \\
\bottomrule
\end{tabular}

\medskip
A is better in \emph{both} subgroups, but B wins in aggregate because 77\% of B's sample is in the easy subgroup, while only 25\% of A's is.
\end{formaldef}

\vspace{8pt}

\begin{intuition}[The Berkeley Admissions Paradox]
In 1973, UC Berkeley's aggregate admission rates showed apparent gender bias: 44\% of men admitted vs.\ 35\% of women. But department-by-department, women were admitted at equal or higher rates. The explanation: women disproportionately applied to competitive departments with low admission rates.

The ``paradox'' vanishes once you realize that aggregation ignores a lurking variable (department competitiveness) whose distribution differs across groups. \textbf{Simpson's Paradox is really about weighted averages with different weights.}
\end{intuition}

\vspace{8pt}

\begin{application}[Dashboard: Simpson's Paradox Detection]
The Segments tab includes a Simpson's Paradox callout. The backend helper \code{\_simpsons\_paradox(df)} in \code{routers/segments.py} computes:
\begin{enumerate}[nosep]
  \item The aggregate treatment effect direction
  \item Per-segment effect directions (by \code{user\_segment})
  \item Paradox flag: True if the aggregate direction differs from \emph{all} segment directions
\end{enumerate}

In the enriched dataset, the treatment works differently for new vs.\ returning users (by design): the treatment helps new users but slightly hurts returning users. This creates a Simpson's Paradox opportunity when the segment proportions are imbalanced.

\medskip
\textbf{Causal interpretation} (Pearl): condition on $Z$ if it is a confounder; do NOT condition if it is a mediator or collider. In the dashboard, device type and user segment are confounders (they exist before treatment assignment).
\end{application}

\vspace{8pt}

\begin{dashlink}
\textbf{Tab:} Segments (Simpson's Paradox callout section).\\
\textbf{Backend:} \code{routers/segments.py}, helper \code{\_simpsons\_paradox(df)}.\\
\textbf{Also:} \code{\_segment\_effects(df, dimension)} runs z-tests per segment across 5 dimensions.
\end{dashlink}


% ==================================================================
\section{Multiple Testing Corrections}
\label{sec:multiple}
% ==================================================================

\begin{formaldef}[Controlling Error Rates Across Multiple Tests]
When testing $m$ hypotheses at level $\alpha$, the probability of at least one false positive under the global null is:
\[
P(\text{at least one false positive}) = 1 - (1-\alpha)^m
\]
For $m = 20$: $P = 1 - 0.95^{20} = 0.642$.

\medskip
\textbf{Family-Wise Error Rate (FWER):} $P(\text{at least one false positive}) \leq \alpha$.

\textbf{False Discovery Rate (FDR):} $E\!\left[\frac{\text{false positives}}{\text{total rejections}}\right] \leq q$.

\medskip
\textbf{Bonferroni:} Reject $H_i$ if $p_i < \alpha/m$. Proof via union bound:
\[
P\!\left(\bigcup_{i \in \mathcal{H}_0}\{p_i < \alpha/m\}\right) \leq \sum_{i \in \mathcal{H}_0} \frac{\alpha}{m} \leq \alpha
\]
Very conservative: with $m=100$, each test uses $\alpha/100 = 0.0005$.

\medskip
\textbf{Holm (step-down, FWER):} Order p-values $p_{(1)} \leq \cdots \leq p_{(m)}$. For $k = 1, \ldots, m$: if $p_{(k)} < \alpha/(m-k+1)$, reject and continue; else stop. Strictly more powerful than Bonferroni because after each rejection, fewer hypotheses remain.

\medskip
\textbf{Benjamini--Hochberg (FDR):} Find the largest $k$ such that $p_{(k)} \leq (k/m) \cdot q$. Reject $H_{(1)}, \ldots, H_{(k)}$. Controls FDR at level $q$ under independence.

\medskip
\begin{tabular}{@{}llll@{}}
\toprule
Method & Controls & Threshold & Conservatism \\
\midrule
Bonferroni & FWER & $\alpha/m$ & Very high \\
Holm & FWER & Step-down from $\alpha/m$ & Moderate \\
BH & FDR & $(k/m) \cdot q$ & Lowest (different metric) \\
\bottomrule
\end{tabular}
\end{formaldef}

\vspace{8pt}

\begin{intuition}[Three Inspectors]
\textbf{Bonferroni} (paranoid): checking 20 products, uses 0.25\% error per product to keep total risk at 5\%. Safe but misses real defects.

\textbf{Holm} (sequential): starts with the most suspicious item. After flagging it, relaxes the threshold slightly, since fewer suspects remain. Strictly better than Bonferroni.

\textbf{BH-FDR} (pragmatic): ``I accept that among my discoveries, up to 5\% might be false. I will verify the important ones later.'' The right mindset for exploratory segment analysis.
\end{intuition}

\vspace{8pt}

\begin{application}[Dashboard: Segment Analysis Context]
The Segments tab runs z-tests across 5 dimensions (device, browser, country, user segment, traffic source), each with multiple levels. This produces many simultaneous tests.

\medskip
\textbf{Practical recommendation for the dashboard:}
\begin{enumerate}[nosep]
  \item \textbf{Primary metric} (conversion rate on the Overview tab): tested without correction. This is the experiment's official success criterion.
  \item \textbf{Secondary metrics} (revenue, session duration on the Frequentist tab): corrected with Holm.
  \item \textbf{Segment breakdowns} (Segments tab): corrected with BH-FDR, since these are exploratory.
  \item \textbf{Guardrail metrics} (if any): separate FWER correction.
\end{enumerate}

\medskip
\textbf{Worked example} (5 metrics tested simultaneously):

\medskip
\begin{tabular}{@{}lcccc@{}}
\toprule
Metric & Raw $p$ & Bonferroni & Holm & BH \\
\midrule
Add-to-cart & 0.003 & 0.015 & 0.015 & 0.015 \\
Revenue/visitor & 0.018 & 0.090 & 0.072 & 0.045 \\
Bounce rate & 0.032 & 0.160 & 0.096 & 0.053 \\
Time on page & 0.085 & 0.425 & 0.170 & 0.106 \\
Scroll depth & 0.220 & 1.000 & 0.220 & 0.220 \\
\bottomrule
\end{tabular}

\medskip
All methods agree: add-to-cart is significant. BH additionally flags revenue. The choice of method determines how aggressive your discovery threshold is.
\end{application}


% ==================================================================
\section{Cross-Topic Connections}
% ==================================================================

Understanding how these concepts link together is crucial for mastering A/B testing as a system, not just a collection of techniques.

\begin{enumerate}
  \item \textbf{Power analysis (Sec.~\ref{sec:power}) depends on effect size (Sec.~2.4):} The sample size formula uses the expected $\delta = p_B - p_A$, which can be expressed via Cohen's $h$. Smaller $h$ means more users needed.

  \item \textbf{Sequential testing (Sec.~\ref{sec:sequential}) modifies power analysis:} OBF designs require ${\sim}3\%$ more users than fixed-sample designs. Factor this into runtime estimation.

  \item \textbf{SRM detection (Sec.~1) uses the same chi-squared test as the analysis (Sec.~2.2):} Both compare observed vs.\ expected frequencies. SRM is a ``meta-test'' that validates the infrastructure before you trust the treatment effect test.

  \item \textbf{Simpson's Paradox (Sec.~\ref{sec:simpson}) motivates segmented analysis, which triggers multiple testing (Sec.~\ref{sec:multiple}):} When you segment results to check for the paradox, you run many tests. Apply BH-FDR to segment-level comparisons.

  \item \textbf{Bayesian expected loss (Sec.~3.3) provides a natural alternative to multiple testing correction:} The loss function inherently penalizes overconfident conclusions, reducing (but not eliminating) the multiple comparisons problem.

  \item \textbf{Frequentist CIs (Sec.~2) converge to Bayesian credible intervals (Sec.~3) with flat priors:} The Bernstein--von Mises theorem (Sec.~4) guarantees this for large $n$. For small $n$, informative priors can be substantially more accurate.

  \item \textbf{Sequential testing and Bayesian stopping rules:} Both address peeking. Frequentist sequential testing controls Type~I error via alpha spending. Bayesian stopping controls expected loss. The Bayesian approach is simpler but lacks the Type~I error guarantee.
\end{enumerate}


% ==================================================================
\section{Summary Table}
% ==================================================================

\begin{tabular}{@{}p{3.5cm}p{5cm}p{2cm}p{3.5cm}@{}}
\toprule
\textbf{Concept} & \textbf{Core Formula} & \textbf{Dashboard Tab} & \textbf{Backend Function} \\
\midrule
Z-test (proportions) & $Z = (\hat{p}_B - \hat{p}_A) / \mathrm{SE}_{\mathrm{pool}}$ & Frequentist & \code{z\_test\_proportions} \\
Chi-squared & $\chi^2 = \sum (O-E)^2/E$ & Frequentist & \code{chi\_squared\_test} \\
Wilson CI & $\tilde{p} \pm z \cdot [\ldots] / (1+z^2/n)$ & Frequentist & \code{wilson\_confidence\_interval} \\
Cohen's $h$ & $2\arcsin(\sqrt{p_B}) - 2\arcsin(\sqrt{p_A})$ & Frequentist & \code{cohens\_h} \\
Beta posterior & $\mathrm{Beta}(\alpha+X, \beta+n-X)$ & Bayesian & \code{beta\_posterior} \\
$P(B>A)$ & $\frac{1}{S}\sum \mathbf{1}[p_B^{(s)} > p_A^{(s)}]$ & Bayesian & \code{probability\_b\_beats\_a} \\
Expected loss & $\frac{1}{S}\sum \max(p_A^{(s)}-p_B^{(s)}, 0)$ & Bayesian & \code{expected\_loss} \\
Sample size & $n = [(z_\alpha\sqrt{\ldots}+z_\beta\sqrt{\ldots})/\delta]^2$ & Power & \code{required\_sample\_size} \\
MDE & Binary search on sample size & Power & \code{minimum\_detectable\_effect} \\
OBF boundary & $c_k = c\sqrt{K/k}$ & Sequential & \code{obrien\_fleming\_bounds} \\
SRM test & $\chi^2 = \sum(O_i-E_i)^2/E_i$ & Overview & \code{sample\_ratio\_mismatch\_test} \\
\bottomrule
\end{tabular}


% ==================================================================
\section{Key Formulas Quick Reference}
% ==================================================================

\begin{align}
\text{Pooled proportion:} \quad & \hat{p} = \frac{X_A + X_B}{n_A + n_B} \\[6pt]
\text{Z-statistic:} \quad & Z = \frac{\hat{p}_B - \hat{p}_A}{\sqrt{\hat{p}(1-\hat{p})(1/n_A + 1/n_B)}} \\[6pt]
\text{Wilson CI center:} \quad & \tilde{p} = \frac{\hat{p} + z^2/(2n)}{1 + z^2/n} \\[6pt]
\text{Cohen's } h: \quad & h = 2\arcsin\!\left(\sqrt{p_B}\right) - 2\arcsin\!\left(\sqrt{p_A}\right) \\[6pt]
\text{Beta posterior:} \quad & \mathrm{Beta}(\alpha + X, \;\beta + n - X) \\[6pt]
\text{P}(B > A): \quad & \frac{1}{S}\sum_{s=1}^S \mathbf{1}\!\left[p_B^{(s)} > p_A^{(s)}\right] \\[6pt]
\text{Expected loss:} \quad & \frac{1}{S}\sum_{s=1}^S \max\!\left(p_A^{(s)} - p_B^{(s)}, 0\right) \\[6pt]
\text{Sample size:} \quad & n = \frac{(z_{\alpha/2} + z_\beta)^2 \cdot [p_A(1{-}p_A) + p_B(1{-}p_B)]}{(p_B - p_A)^2} \\[6pt]
\text{Runtime:} \quad & \text{Days} = \frac{2n}{T \cdot f} \\[6pt]
\text{OBF boundary:} \quad & c_k = c \cdot \sqrt{K/k} \\[6pt]
\text{OBF spending:} \quad & \alpha^*(t) = 2 - 2\Phi\!\left(\frac{z_{\alpha/2}}{\sqrt{t}}\right) \\[6pt]
\text{Bonferroni:} \quad & \text{Reject if } p_i < \alpha / m
\end{align}

\end{document}
